\documentclass[11pt,twoside,a4paper]
{article}

\usepackage{graphicx} 
\usepackage[T1]{fontenc}
\usepackage[polish]{babel}
\usepackage[utf8]{inputenc}
\usepackage{lmodern}
\selectlanguage{polish}
\usepackage{indentfirst}
\usepackage{cite}

\begin{document}

\begin{figure}
\centering
\includegraphics[scale=1]{logoAGH.jpg} 
\end{figure}
\title{\textbf{M-Health}\medskip\\Referat na przedmiot Systemy Informatyczne w Medycynie\\}
\author{Autorzy:\\Agata Paciorek\\Grzegorz Bylina\\Kierunek Informatyka\\Wydział Elektrotechniki Automatyki Informatyki i Inżynierii Biomedycznej\bigskip\\}
\date{Kraków, 2014}
\maketitle
\newpage
\tableofcontents
\newpage
\section{Wstęp teoretyczny}
Koncepcja \textbf{M-Health} (\emph{Mobile Health}) jest integracją mobilnej telekomunikacji i technologii multimedialnej, której celem jest dostarczanie opieki zdrowotnej.\\
Zadania M-Health \cite{5969916} :
\begin{itemize}
\item zbieranie danych klinicznych
\item dostarczanie informacji do personelu medycznego, badaczy, pacjentów
\item monitorowanie pacjenta real-time
\item dostarczanie opieki medycznej
\end{itemize}
Mobile Health zajmuje się dostarczaniem opieki wszędzie i o każdej porze, surpassing geographical, temporal, even organizational barier \cite{6655256}. Nie kończy się (m-health) na udostępnianiu aplikacji opieki zdrowotnej na urządzeniach mobilnych as m-health can involve czujniki i sieci bezprzewodowe, urządzenia mobilne do dostępu do variety of usług medycznych, profesjoalistów w celu podjęcia decyzji i dostarczaniu opieki zdrowotnej i zarządzaniu aktywności dla osób starszych \cite{Varshney2014}.Weź to popraw człowieku zeby nie bylo żywcem z artykulu.

\section{Przykłady rozwiązań}
\subsection{Monitorowanie pacjenta}
\paragraph{Monitorowanie postury \cite{Lee:2012:MPM:2370216.2370320}\\}
Zaproponowany mechanizm szacuje różne wartości reprezentujące posturę użytkownika takie jak: kąt nachylenia szyi , odległość oglądania, stan wzroku użytkownika poprzez analizę danych z przedniej kamery , akcelerometru, czujnika orientacji lub dowolną ich kombinację. Powiadamia użytkownika jeśli oszacowane
wartości są utrzymywane w nieprawidłowym zakresie ponad dozwolony czas.
\paragraph{Ocena zdrowia i zdolności kierowcy \cite{6583819}\\}
Aplikacja sprzężona z Data Loggerem zainstalowanym w pojeździe, który gromadzi dane z czujników znajdujących się w pojeździe i na ciele kierowcy. Aplikacja jest przeznaczona do monitorowania jakości sygnałów real-time i znacznego skrócenia czasu oszacowania stanu zdrowia kierowcy łącznie z jego zdolnościami do prowadzenia pojazdów. System składa się z dwóch modułów:
\begin{itemize}
\item smartbio 1 - system badania\\
Służy do ułatwienia gromadzenia informacji na temat rożnych badań fizjologicznych/psychologicznych i wykorzystanie ich do obliczenia współczynnika, określającego zdolność kierowcy do jazdy

\item smartbio 2 - system monitorowania\\
Dzieli się na dwie warstwy hardware i monitoring. Warstwa hardware gromadzi i wstępnie przetwarza dane (z sensorów). Warstwa monitoring ma postać aplikacji na smartfona, jej głównym zadaniem jest komunikacja między użytkownikiem, a warstwą hardware
\end{itemize}
\paragraph{Diagnozowanie depresji\\}
 Ta praca prponuje techniki poprawiające komunikacje w BSN (Body Sensor Network), które gromadzą dane o stanach emocjonalnych pacjenta. 
 Te BSN  mogą stale monitorować , dyskretnie szacować i klasyfikować każdy stan depresyjny. Dodatkowo dane na temat życia pacjenta mogą zostać skorelowane z 
 uwarunkowaniami fizjologiczymi aby zidentyfikować jak poszczególne bodźce wywołują objawy. Taki ciągły strumień danych jest poprawą w stosunku do 
 migawki objawów , które zaobserwuje lekarz w ciągu badania. 
 
 
\subsection{Zarządzanie informacją}
\paragraph{Zarządzanie chorobami przewlekłymi\\}
Jest to pewna znaleziona koncepcja aplikacji mobilnej. Celem pracy było zidentyfikowanie funkcji i wymagań funkcyjnych , które pomogły by użytkownikowi w zarządzaniu opieką nad swoją chorobą przewlekłą
 Projekt przewiduje kompleksowe działania takie jak wyświetlanie i zarządzanie lekami, harmonogram badań/wizyt, notatki, plany diet, ważne informacje o planie leczenia

\paragraph{Zarządzanie lekami\\}
 Aplikacja do zarządzania podawaniem leków. To pamiętnik śledzący i zarządzający lekami w celu zapobiegania błędom medycznym. 
 Poprzez wizje,dźwięk,wibracje  SapoMed przypomina użytkowników o ich harmonogramie leków
 Aplikacja umożliwia rejestrowanie leków poprzez kamerę za pomocą barcode na opakowaniu. 
 Korzysta z usług internetowych (web services) aby uzyskać informacje o lekach a nawet o ich dawkowaniu
 Wykorzystuje te web services również do zapamiętywania spożywanych wczesniej leków 
\subsection{Problem bezpieczeństwa}
\section{Zastosowania}
\section{Podsumowanie}

\newpage
\bibliographystyle{alpha}
\bibliography{cytowaniaBylina}
\end{document}